%%%%%%%%%%%%%%%%%%%%%%%%
%
% $Autor: Gautam Ramesh$
% $Datum: 2025-10-17 11:27:47Z $
% $Pfad: /Users/gautamramesh/githubNew/ML26-03-Car-Licence-Plate-Identification-with-a-Jetson-Nano/report/System/EdgeComputer/jetsonNano.tex $
% $Version: 1 $
%
% !TeX encoding = utf8
% !TeX root = Rename
%
%%%%%%%%%%%%%%%%%%%%%%%%

\chapter{Hardware Verification and Testing} \label{ch:hardware_test}

% --- Custom Listing Style for Boxed Code with Doxygen Comments ---
\lstset{
	language=Python,
	basicstyle=\ttfamily\footnotesize,
	backgroundcolor=\color{gray!5},   % Light gray background
	frame=single,                     % Box around the code
	rulecolor=\color{black},          % Border color
	keywordstyle=\color{blue}\bfseries,
	commentstyle=\color{green!40!black}\itshape, % Green for comments
	stringstyle=\color{orange},
	numbers=left,                     % Line numbers on the left
	numberstyle=\tiny\color{gray},
	stepnumber=1,
	numbersep=10pt,
	tabsize=4,
	showspaces=false,
	showstringspaces=false,
	breaklines=true,
	breakatwhitespace=true,
	captionpos=b,                     % Caption at bottom
	escapeinside={(*@}{@*)},          % For LaTeX inside code if needed
	morecomment=[l][\color{magenta}]{\#} % Special color for Doxygen tags if using #
}

\section{Introduction}

An important prerequisite to deploying Deep Learning Models such as YOLOv8 and Optical Character Recognition (OCR) is making sure the physical layer is reliable. 
In the context of this project, the hardware architecture comprises of an edge computing unit along with an optical acquisition unit. Before software pipeline deployment, these components should be verified systematically - especially the Seeed Studio reComputer J1010 and the Microsoft LifeCam Camera interface. 
 
This chapter is a documentation of utilized hardware components, defining their specific functions within the License Plate Recognition (LPR) architecture, while providing an understanding around the rigorous testing protocols applied for operational status verification.

\section{Hardware Components and Specification}

The system builds upon two hardware pillars: the processing unit (Jetson Nano) and the image acquisition unit (Microsoft LifeCam).

\subsection{Processing Unit: reComputer J1010}
The core processing unit used is the \textbf{Seeed Studio reComputer J1010}. The device is an embedded edge AI box consisting the NVIDIA Jetson Nano module (4GB B01). It is selected primarily due to its ability to perform CUDA computations in parallel, which is an important prerequisite for real time YOLO inference. This enables the system to follow a low power profile.  

\begin{table}[h]
	\centering
	\caption{Technical Specifications of reComputer J1010 (Jetson Nano) \cite{Stu23a}}
	\label{tab:jetson_specs}
	\begin{tabular}{@{}llp{6.5cm}@{}}
		\toprule
		\textbf{Component} & \textbf{Specification} & \textbf{Function in LPR System} \\ \midrule
		\textbf{Module} & Jetson Nano 4GB & The brain of the system running Ubuntu Linux (L4T). \\
		\textbf{GPU} & 128-core Maxwell & Accelerates tensor operations for YOLOv8 inference. \\
		\textbf{CPU} & Quad-core ARM A57 & Handles OS tasks, image preprocessing, and database logging. \\
		\textbf{Memory} & 4 GB 64-bit LPDDR4 & Stores loaded neural network weights and frame buffers. \\
		\textbf{Storage} & 16 GB eMMC 5.1 & Hosts the operating system and Python libraries. \\
		\textbf{Interface} & USB 3.0 / 2.0 & Provides high-speed connectivity for the camera. \\ 
		\bottomrule
	\end{tabular}
\end{table}

\subsection{Optical Unit: Microsoft LifeCam}
For image acquisition, the \textbf{Microsoft LifeCam} (specifically the LifeCam Studio/Show models as referenced in the system design) was selected. This is a high-definition USB webcam capable of capturing the detailed frames necessary for accurate character segmentation during the OCR phase.

\begin{itemize}
	\item \textbf{Sensor Resolution:} 2.0 Megapixel (1080p / 720p video modes).
	\item \textbf{Lens Type:} Wide-angle glass element lens with high-precision optics.
	\item \textbf{Interface:} USB 2.0 (Universal Serial Bus) High Speed.
	\item \textbf{Focus:} Auto-focus capability (critical for focusing on moving vehicles at varying depths).
	\item \textbf{Mounting:} Flexible attachment base for positioning on tripods or monitors.
\end{itemize}

\section{Hardware Functions}

Each component plays a distinct, interconnected role in the LPR pipeline. The data flow between these hardware parts is critical for system latency and accuracy.

\begin{enumerate}
	\item \textbf{Image Acquisition (Camera):} The Microsoft LifeCam is able to captures raw video frames of vehicles as they approach. Primarily, the hardware function is to generate and buffer video frames. Subsequently these frames are transmitted through the USB bus and delivered to the Jetson Nano's memory controller.
	\item \textbf{Preprocessing (CPU):} The CPU - ARM Cortex A57, built on the Jetson Nano is responsible for initial resizing and subsequently normalizes the image frame (e.g., resizing $1920 \times 1080$ to $640 \times 640$) so that it matches the input tensor size mandated by the YOLO network.
	\item \textbf{Inference (GPU):} The 128-core Maxwell GPU is able to process and execute the forward pass of the YOLO convolutional neural network. This constitutes one of the most computationally intensive functions as it requires CUDA cores to perform computationally expensive matrix multiplications in parallel to be able to identify license plate bounding boxes.
\end{enumerate}

\begin{figure}[h]
	\centering
		\begin{tikzpicture}[node distance=3cm, auto,
		block/.style={rectangle, draw=blue!60, fill=blue!5, very thick, minimum size=1cm, minimum width=3cm, rounded corners},
		line/.style={draw, -latex', thick, blue}]
		
		% Nodes
		\node [block] (cam) {Microsoft LifeCam};
		\node [block, right=of cam, xshift=1cm] (jetson) {reComputer J1010};
		\node [block, below=of jetson] (display) {Monitor (HDMI)};
		
		% Text labels inside nodes
		\node[below=0.1cm of cam, font=\footnotesize] {Acquisition};
		\node[below=0.1cm of jetson, font=\footnotesize] {Processing (CPU+GPU)};
		\node[below=0.1cm of display, font=\footnotesize] {Visualization};
		
		% Edges
		\path [line] (cam) -- node[midway, above] {USB 2.0 (Raw Frames)} (jetson);
		\path [line] (jetson) -- node[midway, right] {HDMI Output} (display);
		
	\end{tikzpicture}
	\caption{Hardware Data Flow Topology}
	\label{fig:hw_topology}
\end{figure}

\section{Hardware Testing Methodology}

To ensure the hardware supports the software requirements described in previous chapters, we devised specific test cases. The objective was to validate the operational integrity of the GPU and the Camera interface using Python scripts.

\subsection{Test 1: GPU and Environment Verification}
Since the YOLOv8 model relies heavily on GPU acceleration to achieve real-time frame rates (FPS), verifying the GPU availability is the first step. A Python script leveraging the PyTorch library to query the hardware status is utilized. The test validates and provides confirmation to the JetPack OS as to whether it is correctly installed and that the CUDA libraries are accessible within the Python environment.

\textbf{Test Procedure:}
\begin{enumerate}
	\item Power and boot the reComputer J1010.
	\item Navigate to the testing directory by opening a terminal window.
	\item Perform successful verification by executing the Python Verification script.
\end{enumerate}

The following code, \texttt{test\_gpu.py}, was executed on the system:

\begin{lstlisting}[language=Python, caption={Python Script for CUDA and GPU Verification}, label={lst:gpu_test}]
	## @file test_gpu.py
	#  @brief Verifies CUDA and GPU availability for YOLOv8 on Jetson Nano.
	#  @details This script checks the PyTorch version, verifies that the CUDA
	#           drivers are accessible, and attempts to allocate a tensor in
	#           GPU memory to confirm hardware readiness.
	#  @author [Your Name]
	#  @date 2025-10-25
	
	import torch
	import sys
	
	## @brief Main verification function.
	def verify_system():
	print("--- Hardware Verification: Jetson Nano ---")
	
	# 1. Check Python Version
	print(f"Python Version: {sys.version}")
	
	# 2. Check CUDA Availability for YOLO
	# This checks if the NVIDIA Drivers are loaded
	cuda_available = torch.cuda.is_available()
	print(f"CUDA Available: {cuda_available}")
	
	if cuda_available:
	# 3. Get Device Properties
	# This queries the Maxwell GPU specifically
	device = torch.device("cuda")
	props = torch.cuda.get_device_properties(device)
	print(f"Device Name: {props.name}")
	print(f"Total Memory: {props.total_memory / 1e9:.2f} GB")
	print(f"Multi-Processor Count: {props.multi_processor_count}")
	
	# 4. Perform a Tensor Calculation Test
	# This verifies write/read access to GPU memory
	try:
	x = torch.rand(5, 3).to(device)
	print("Tensor Allocation Test: SUCCESS")
	print(f"Test Tensor: \n{x}")
	except Exception as e:
	print(f"Tensor Allocation Test: FAILED - {e}")
	else:
	print("CRITICAL WARNING: System is running on CPU only.")
	
	if __name__ == "__main__":
	verify_system()
\end{lstlisting}

\textbf{Results of Test:} \\
The output confirms correctness in terms of configuration of the reComputer J1010. The system was able to identify the "NVIDIA Tegra X1" device and performed successful test tensor allocation. Memory checks reported roughly 3.9 GB of accessible memory (shared with the CPU), confirming hardware readiness for YOLO training and inference phases.

\subsection{Test 2: Camera Interface and Real-Time Capture}

The camera interface was tested by employing the standard Video4Linux2 (V4L2) backend via OpenCV. The primary goal was to ascertain whether camera was able to initialize correctly, capture frames, and release the resource without generating errors. This test uses code logic described in the Hardware Description (Chapter 4 of the report).

\textbf{Test Logic Flow:}
Figure \ref{fig:cam_test_flow} illustrates the logic used to verify the camera hardware.

\begin{figure}[h]
	\centering
		\begin{tikzpicture}[node distance=1.5cm, auto,
		startstop/.style={rectangle, rounded corners, minimum width=3cm, minimum height=1cm,text centered, draw=black, fill=red!30},
		process/.style={rectangle, minimum width=3cm, minimum height=1cm, text centered, draw=black, fill=orange!30},
		decision/.style={diamond, minimum width=3cm, minimum height=1cm, text centered, draw=black, fill=green!30},
		arrow/.style={thick,->,>=stealth}]
		
		\node (start) [startstop] {Start Test};
		\node (init) [process, below of=start] {Init VideoCapture(0)};
		\node (check) [decision, below of=init, yshift=-0.5cm] {Is Opened?};
		\node (fail) [process, right of=check, xshift=3cm] {Print Error};
		\node (read) [process, below of=check, yshift=-1cm] {Read Frame};
		\node (display) [process, below of=read] {Display Window};
		\node (stop) [startstop, below of=display] {Release \& Exit};
		
		\draw [arrow] (start) -- (init);
		\draw [arrow] (init) -- (check);
		\draw [arrow] (check) -- node[anchor=south] {No} (fail);
		\draw [arrow] (check) -- node[anchor=east] {Yes} (read);
		\draw [arrow] (read) -- (display);
		\draw [arrow] (display) -- (stop);
		\draw [arrow] (fail) |- (stop);
	\end{tikzpicture}
	\caption{Logic Flow for Camera Hardware Verification}
	\label{fig:cam_test_flow}
\end{figure}

The test code, adapted from the \texttt{LifeCam.py} example provided in the system documentation, is shown below. It initializes the camera at index 0 and attempts to stream video.

\begin{lstlisting}[language=Python, caption={OpenCV Script for Camera Hardware Testing (LifeCam.py)}, label={lst:cam_test}]
	## @file test_camera.py
	#  @brief Tests the Microsoft LifeCam hardware interface.
	#  @details This script initializes the default camera (index 0), performs
	#           a warm-up sequence, and captures a live video stream to verify
	#           sensor functionality and USB bandwidth.
	#  @param source The camera index to open (default: 0).
	
	import cv2
	import time
	
	## @brief Initializes camera and displays video feed.
	#  @param source Index of the camera device.
	def test_camera_hardware(source=0):
	print(f"--- Testing Camera Source: {source} ---")
	
	# 1. Initialize VideoCapture
	# Source 0 is typically the USB camera on Jetson Nano
	cap = cv2.VideoCapture(source)
	
	# Set resolution to standard VGA or 720p for testing
	cap.set(cv2.CAP_PROP_FRAME_WIDTH, 640)
	cap.set(cv2.CAP_PROP_FRAME_HEIGHT, 480)
	
	if not cap.isOpened():
	print("ERROR: Could not open camera hardware.")
	return
	
	# 2. Warm-up camera (Auto-exposure/Auto-focus adjustment)
	# This allows the LifeCam sensor to adjust to ambient light
	print("Warming up sensor...")
	time.sleep(2)
	
	# 3. Capture a Loop of Frames (Simulation of Live Feed)
	print("Starting video stream. Press 'q' to exit.")
	while True:
	ret, frame = cap.read()
	
	if ret:
	# Display the frame to verify visual clarity
	cv2.imshow("Hardware Test: LifeCam", frame)
	
	# Press 'q' to break the loop
	if cv2.waitKey(1) & 0xFF == ord('q'):
	break
	else:
	print("Frame Capture: FAILED")
	break
	
	# 4. Release Hardware Resource
	cap.release()
	cv2.destroyAllWindows()
	print("Camera resource released.")
	
	if __name__ == "__main__":
	test_camera_hardware(0)
\end{lstlisting}

\textbf{Test Results:} \\
During execution runs, the script initializes the Microsoft LifeCam with success. There exists a period known as the "warm-up" that allows the camera's auto-focus mechanism to stabilize. The video feed, displayed in an OpenCV window, verified that the USB bandwidth on the reComputer J1010 was sufficient to display a video stream without frame loss. Subsequently, image clarity was checked visually to ensure it met requirements for OCR processing steps that follow.

\section{Hardware Tests Conclusion}

Hardware verification phase confirmed that the reComputer J1010 and Microsoft LifeCam are interoperable and compatible. The GPU is accessible for deep learning tasks and the camera is able to  provide consistent and visually stable video streams with the use of OpenCV. Such tests validate hardware foundations and serve as an important requirement for the following chapters on YOLO detection and optical character recognition.